\chapter{Einleitung}
\label{cha:einleitung}

% Die Einleitung führt vom Allgemeinen zum Konkreten: Warum ist das Thema
% relevant? Welches Problem wird gelöst? Was genau ist das Ziel? Wie ist die
% Arbeit aufgebaut? Richtwert: 5–10 % des Umfangs.

\section{Motivation}
\label{sec:motivation}

% Kontext und Relevanz des Themas – gern mit einer aktuellen Quelle belegen.
Wissenschaftliche Arbeiten werden häufig mit \LaTeX{} gesetzt, weil Gliederung,
Querverweise und Literaturverwaltung zuverlässig funktionieren~\cite{lamport1994latex}.
Gleichzeitig \ldots

\section{Problemstellung}
\label{sec:problemstellung}

% Welche konkrete Lücke oder Schwierigkeit besteht? Was passiert, wenn sie
% nicht gelöst wird?
Bisherige Ansätze \ldots{} Daraus ergibt sich das Problem, dass \ldots

\section{Zielsetzung und Forschungsfrage}
\label{sec:zielsetzung}

% Ein klar überprüfbares Ziel und eine (bis drei) Forschungsfragen. Diese
% Fragen werden in Kapitel~\ref{cha:fazit} explizit beantwortet.
Ziel dieser Arbeit ist es, \ldots{} Daraus leitet sich folgende Forschungsfrage ab:

\begin{quote}
  \textbf{FF1:} Inwiefern \ldots?
\end{quote}

Zur Beantwortung werden folgende Teilziele verfolgt:
\begin{enumerate}
  \item Analyse der Anforderungen an \ldots
  \item Konzeption und prototypische Umsetzung von \ldots
  \item Evaluation anhand von \ldots
\end{enumerate}

\section{Abgrenzung}
\label{sec:abgrenzung}

% Was wird bewusst nicht behandelt? Das schützt vor der Frage "Warum fehlt X?".
Nicht Gegenstand dieser Arbeit sind \ldots

\section{Aufbau der Arbeit}
\label{sec:aufbau}

Kapitel~\ref{cha:grundlagen} erläutert die theoretischen Grundlagen und ordnet
verwandte Arbeiten ein. Kapitel~\ref{cha:methodik} beschreibt das methodische
Vorgehen und die Anforderungen. Kapitel~\ref{cha:umsetzung} dokumentiert Konzept und
Implementierung. Kapitel~\ref{cha:evaluation} bewertet die Ergebnisse, bevor
Kapitel~\ref{cha:fazit} die Arbeit zusammenfasst und einen Ausblick gibt.
